% !Mode:: "TeX:UTF-8"
% !TEX root = ../../main.tex
% ============================================================
% TOPIC 8
% ============================================================
\mashgulot{8}{The matrix method for systems of linear equations}

\begin{tarif}
Let $Ax = b$ be a system of $n$ linear equations in $n$ unknowns. If the
coefficient matrix $A$ is nonsingular ($\det A \neq 0$), the system has a
unique solution, given by
\[\kf{x = A^{-1}\cdot b,}\]
where $A^{-1}$ is the inverse of $A$ and $b$ is the column vector of
constant terms.
\end{tarif}

The method follows directly from the definition. We write down the
coefficient matrix $A$ and the column $b$ and compute $\det A$; if
$\det A = 0$, the matrix method cannot be applied. Otherwise we find
$A^{-1}$, either from the cofactors or by Gauss–Jordan elimination, and
form the product $x = A^{-1}b$. Substituting the vector $x$ back into the
original equations gives a quick check of the arithmetic.

We stress that the matrix method applies only to square systems
($n$ equations in $n$ unknowns) with $\det A \neq 0$, that is, to systems
with exactly one solution.

\subsection{Solving $2\times 2$ and $3\times 3$ systems by the matrix method}

\textbf{The $2\times 2$ case.}

\begin{misol}
$\begin{cases}3x_1 + x_2 = 5,\\ 5x_1 + 2x_2 = 7.\end{cases}$

\textbf{Solution.}
Here
$A = \begin{pmatrix}3&1\\5&2\end{pmatrix}$ and
$b = \begin{pmatrix}5\\7\end{pmatrix}$.
Since $\det A = 3\cdot 2 - 1\cdot 5 = 1 \neq 0$, the system has a unique
solution. For a $2\times 2$ matrix the inverse is obtained by swapping the
diagonal entries, changing the signs of the off-diagonal ones and dividing
by the determinant:
\[\kf{A^{-1} = \frac{1}{1}\begin{pmatrix}2&-1\\-5&3\end{pmatrix}
= \begin{pmatrix}2&-1\\-5&3\end{pmatrix}.}\]
Multiplying $A^{-1}$ by $b$, we obtain
\[\kf{x = A^{-1}b = \begin{pmatrix}2&-1\\-5&3\end{pmatrix}\begin{pmatrix}5\\7\end{pmatrix}
= \begin{pmatrix}2\cdot 5 + (-1)\cdot 7\\ (-5)\cdot 5 + 3\cdot 7\end{pmatrix}
= \begin{pmatrix}10-7\\-25+21\end{pmatrix}
= \begin{pmatrix}3\\-4\end{pmatrix}.}\]

\textbf{Answer:} $x_1 = 3$, $x_2 = -4$.

These values satisfy both equations:
$3\cdot 3 + (-4) = 5$ and $5\cdot 3 + 2\cdot(-4) = 7$.
\end{misol}

\begin{misol}
$\begin{cases}-2x_1 + 3x_2 = 1,\\ x_1 - x_2 = 4.\end{cases}$

\textbf{Solution.}
We have
$A = \begin{pmatrix}-2&3\\1&-1\end{pmatrix}$,
$b = \begin{pmatrix}1\\4\end{pmatrix}$ and
$\det A = (-2)(-1) - 3\cdot 1 = -1 \neq 0$. Hence
\[\kf{A^{-1} = \frac{1}{-1}\begin{pmatrix}-1&-3\\-1&-2\end{pmatrix}
= \begin{pmatrix}1&3\\1&2\end{pmatrix},}\]
and
\[\kf{x = \begin{pmatrix}1&3\\1&2\end{pmatrix}\begin{pmatrix}1\\4\end{pmatrix}
= \begin{pmatrix}1+12\\1+8\end{pmatrix}
= \begin{pmatrix}13\\9\end{pmatrix}.}\]

\textbf{Answer:} $x_1 = 13$, $x_2 = 9$.

Check: $-2\cdot 13 + 3\cdot 9 = -26+27 = 1$ and $13-9 = 4$.
\end{misol}

\textbf{The $3\times 3$ case.}

\begin{misol}
$\begin{cases}
x_1 - x_2 + x_3 = 0,\\
2x_2 - x_3 = 1,\\
2x_1 + 3x_2 = 5.
\end{cases}$

\textbf{Solution.}
The coefficient matrix and the column of constants are
$A = \begin{pmatrix}1&-1&1\\0&2&-1\\2&3&0\end{pmatrix}$ and
$b = \begin{pmatrix}0\\1\\5\end{pmatrix}$.
Expanding $\det A$ along the first row, we get
\[\kf{\det A = 1\begin{vmatrix}2&-1\\3&0\end{vmatrix}
-(-1)\begin{vmatrix}0&-1\\2&0\end{vmatrix}
+1\begin{vmatrix}0&2\\2&3\end{vmatrix}
= 1\cdot 3 + 1\cdot 2 + 1\cdot(-4) = 1,}\]
so $A$ is nonsingular. Next we compute the cofactors:
{\small
\begin{align*}
A_{11} &= \begin{vmatrix}2&-1\\3&0\end{vmatrix}=3, &
A_{12} &= -\begin{vmatrix}0&-1\\2&0\end{vmatrix}=-2, &
A_{13} &= \begin{vmatrix}0&2\\2&3\end{vmatrix}=-4, \\[4pt]
A_{21} &= -\begin{vmatrix}-1&1\\3&0\end{vmatrix}=3, &
A_{22} &= \begin{vmatrix}1&1\\2&0\end{vmatrix}=-2, &
A_{23} &= -\begin{vmatrix}1&-1\\2&3\end{vmatrix}=-5, \\[4pt]
A_{31} &= \begin{vmatrix}-1&1\\2&-1\end{vmatrix}=-1, &
A_{32} &= -\begin{vmatrix}1&1\\0&-1\end{vmatrix}=1, &
A_{33} &= \begin{vmatrix}1&-1\\0&2\end{vmatrix}=2.
\end{align*}}
Placing the cofactors of the $i$th row of $A$ in the $i$th column, we obtain
the adjugate matrix, and therefore
\[\kf{A^{-1} = \frac{1}{1}\begin{pmatrix}3&3&-1\\-2&-2&1\\-4&-5&2\end{pmatrix}
= \begin{pmatrix}3&3&-1\\-2&-2&1\\-4&-5&2\end{pmatrix}.}\]
It remains to multiply by $b$:
\[\kf{x = A^{-1}b = \begin{pmatrix}3&3&-1\\-2&-2&1\\-4&-5&2\end{pmatrix}
\begin{pmatrix}0\\1\\5\end{pmatrix}
= \begin{pmatrix}0+3-5\\0-2+5\\0-5+10\end{pmatrix}
= \begin{pmatrix}-2\\3\\5\end{pmatrix}.}\]

\textbf{Answer:} $x_1=-2$, $x_2=3$, $x_3=5$.

Substitution confirms the result:
$(-2)-3+5 = 0$, \ $2\cdot 3 - 5 = 1$, \ $2\cdot(-2)+3\cdot 3 = 5$.
\end{misol}

\begin{misol}
$\begin{cases}
2x_1 - x_2 = 1,\\
x_1 - x_3 = -1,\\
3x_1 + 2x_2 + x_3 = 4.
\end{cases}$

\textbf{Solution.}
Here $A = \begin{pmatrix}2&-1&0\\1&0&-1\\3&2&1\end{pmatrix}$,
$b = \begin{pmatrix}1\\-1\\4\end{pmatrix}$, and expansion along the first
row gives $\det A = 2(0+2)+1\cdot(1+3)+0 = 8\neq 0$. The cofactors are
\begin{align*}
A_{11}&=2,& A_{12}&=-4,& A_{13}&=2,\\
A_{21}&=1,& A_{22}&=2,& A_{23}&=-7,\\
A_{31}&=1,& A_{32}&=2,& A_{33}&=1,
\end{align*}
so that
\[\kf{A^{-1} = \frac{1}{8}\begin{pmatrix}2&1&1\\-4&2&2\\2&-7&1\end{pmatrix}}\]
and
\[\kf{x = \frac{1}{8}\begin{pmatrix}2&1&1\\-4&2&2\\2&-7&1\end{pmatrix}
\begin{pmatrix}1\\-1\\4\end{pmatrix}
= \frac{1}{8}\begin{pmatrix}2-1+4\\-4-2+8\\2+7+4\end{pmatrix}
= \frac{1}{8}\begin{pmatrix}5\\2\\13\end{pmatrix}
= \begin{pmatrix}5/8\\1/4\\13/8\end{pmatrix}.}\]

For instance, the first equation holds:
$2\cdot\frac{5}{8}-\frac{1}{4}=\frac{10}{8}-\frac{2}{8}=1$.
\end{misol}

\textbf{Comparison of the matrix method with Cramer’s rule.}

\begin{center}
\begin{tabular}{|p{5.3cm}|p{4.3cm}|}
\hline
\textbf{Matrix method ($x=A^{-1}b$)} & \textbf{Cramer’s rule ($x_j=\Delta_j/\Delta$)} \\
\hline
Requires the inverse matrix &
Requires only determinants \\
\hline
Once $A^{-1}$ is known, each new $b$ costs one multiplication &
The determinants must be recomputed for every $b$ \\
\hline
$A^{-1}$ has $n^2$ entries &
$n+1$ determinants of order $n$ \\
\hline
Computing $A^{-1}$ is expensive for large $n$ &
The work grows very fast for $n>3$ \\
\hline
Also suits matrix equations such as $AXB=C$ &
Only for systems of the form $Ax=b$ \\
\hline
\end{tabular}
\end{center}

Thus, when several systems with the same matrix $A$ and different right-hand
sides have to be solved, the matrix method is preferable, since $A^{-1}$ is
found only once. For a single small system Cramer’s rule may be simpler.

\subsection{Applications of the matrix method in information security}

\textbf{Encryption and decryption with the Hill cipher.}
The Hill cipher is a direct application of the matrix method. A plaintext
vector $p$ is encrypted by multiplying it by the key matrix $K$, which gives
the ciphertext $c = Kp$. To decrypt, one multiplies the ciphertext by the
inverse matrix: $p = K^{-1}c$, which is exactly the formula $x=A^{-1}b$.
Since every ciphertext block is decrypted with the same matrix $K^{-1}$, the
cipher illustrates well the main advantage of the method: the inverse is
computed once and then used for many multiplications.

\begin{misol}
Let the key matrix be $K=\begin{pmatrix}3&1\\5&2\end{pmatrix}$ and let three
ciphertext blocks $c_1=\begin{pmatrix}5\\7\end{pmatrix}$,
$c_2=\begin{pmatrix}10\\18\end{pmatrix}$, $c_3=\begin{pmatrix}3\\4\end{pmatrix}$
be given.

Having found $K^{-1}=\begin{pmatrix}2&-1\\-5&3\end{pmatrix}$ once, we obtain
\[\kf{p_1 = K^{-1}c_1 = \begin{pmatrix}3\\-4\end{pmatrix},\quad
p_2 = K^{-1}c_2 = \begin{pmatrix}2\\4\end{pmatrix},\quad
p_3 = K^{-1}c_3 = \begin{pmatrix}2\\-3\end{pmatrix}.}\]
All three blocks are decrypted with the same matrix $K^{-1}$.
\end{misol}

\textbf{Linear cryptanalysis.}
In cryptanalysis (the study of breaking ciphers), when several
plaintext–ciphertext pairs are known, finding the key matrix reduces to a
matrix equation $AX=B$ with the unknown key $X$. Its solution is
$X = A^{-1}B$.

\textbf{Decoding in coding theory.}
For systematic codes the matrix method is used to extract the information
part of a received word. If the generator matrix has the form
$G=(I_k\,|\,P)$, the information bits can be read off directly from the
received vector $c$. If $G$ is not in systematic form, one has to apply the
inverse of the information part of $G$.

\textbf{Digital signal processing.}
Linear filters used to filter and analyse signals in information security
(for example, in intrusion detection systems) are described by relations of
the form $y = Hx$. The signal $x$ is recovered by the formula
$x = H^{-1}y$, which is one more instance of the matrix method.

\subsection{Computing with Python, Maple and Mathcad}

Let us solve the system of the third example once more by computer. In
Python the SymPy library computes the determinant, the inverse matrix and
the product $A^{-1}b$ in exact rational arithmetic.
\begin{Verbatim}[fontsize=\footnotesize,frame=single,label={\small Python}]
from sympy import Matrix
A = Matrix([[1, -1, 1], [0, 2, -1], [2, 3, 0]])
b = Matrix([0, 1, 5])
print(A.det())
Ainv = A.inv()
print(Ainv)
print((Ainv*b).T)
\end{Verbatim}
\begin{Verbatim}[fontsize=\footnotesize,frame=single,label={\small Output}]
1
Matrix([[3, 3, -1], [-2, -2, 1], [-4, -5, 2]])
Matrix([[-2, 3, 5]])
\end{Verbatim}

In Maple the same computation uses the \texttt{LinearAlgebra} package; the
command \texttt{LinearSolve(A, b)} gives the solution directly and serves as
a check.
\begin{Verbatim}[fontsize=\footnotesize,frame=single,label={\small Maple}]
with(LinearAlgebra):
A := Matrix([[1, -1, 1], [0, 2, -1], [2, 3, 0]]):
b := Vector([0, 1, 5]):
Determinant(A);           # 1
Ainv := MatrixInverse(A);
x := Ainv . b;            # <-2, 3, 5>
LinearSolve(A, b);        # the same vector
\end{Verbatim}

In Mathcad Prime the matrices are entered through the matrix template, the
inverse is obtained with the exponent $-1$, and the evaluation sign
\texttt{=} displays the result; the built-in function \texttt{lsolve(A, b)}
returns the same vector.
\begin{Verbatim}[fontsize=\footnotesize,frame=single,label={\small Mathcad}]
     | 1  -1   1 |          | 0 |
A := | 0   2  -1 |     b := | 1 |
     | 2   3   0 |          | 5 |

|A| = 1

        |  3   3  -1 |              | -2 |
A^-1 =  | -2  -2   1 |   A^-1*b =   |  3 |
        | -4  -5   2 |              |  5 |

                | -2 |
lsolve(A, b) =  |  3 |
                |  5 |
\end{Verbatim}

\subsection{Exercises on the matrix method}
\bigskip

\emph{1. Solve the following $2\times 2$ systems by the matrix method.}
\bigskip
\begin{multicols}{2}\raggedright
\begin{enumerate}[label=\arabic*., leftmargin=*]

\item $\begin{cases}3x_1+x_2=5\\5x_1+2x_2=7\end{cases}$

\item $\begin{cases}-2x_1+3x_2=1\\x_1-x_2=4\end{cases}$

\item $\begin{cases}x_1+2x_2=8\\3x_1-x_2=3\end{cases}$

\item $\begin{cases}4x_1-x_2=6\\2x_1+3x_2=8\end{cases}$

\item $\begin{cases}-x_1+5x_2=9\\2x_1-3x_2=-1\end{cases}$

\item $\begin{cases}7x_1+2x_2=11\\3x_1+x_2=5\end{cases}$

\item $\begin{cases}x_1-4x_2=-7\\-2x_1+7x_2=11\end{cases}$

\item $\begin{cases}5x_1+3x_2=19\\2x_1-x_2=5\end{cases}$

\item $\begin{cases}-3x_1+2x_2=7\\x_1+x_2=6\end{cases}$

\item $\begin{cases}6x_1-5x_2=1\\-2x_1+3x_2=5\end{cases}$

\item $\begin{cases}x_1+x_2=4\\2x_1-x_2=5\end{cases}$

\item $\begin{cases}8x_1+5x_2=13\\3x_1+2x_2=5\end{cases}$

\item $\begin{cases}-x_1+2x_2=3\\4x_1-5x_2=-2\end{cases}$

\item $\begin{cases}2x_1+7x_2=16\\x_1-3x_2=-5\end{cases}$

\item $\begin{cases}3x_1-4x_2=-1\\x_1+2x_2=7\end{cases}$

\item $\begin{cases}9x_1-x_2=8\\4x_1+x_2=5\end{cases}$

\item $\begin{cases}x_1+6x_2=13\\2x_1-x_2=0\end{cases}$

\item $\begin{cases}5x_1-x_2=4\\-3x_1+2x_2=1\end{cases}$

\item $\begin{cases}4x_1+3x_2=18\\-x_1+2x_2=1\end{cases}$

\item $\begin{cases}x_1-x_2=0\\3x_1+4x_2=7\end{cases}$

\item $\begin{cases}2x_1+x_2=8\\x_1+3x_2=14\end{cases}$

\item $\begin{cases}-4x_1+3x_2=5\\2x_1+x_2=9\end{cases}$

\item $\begin{cases}x_1-5x_2=-9\\3x_1+2x_2=4\end{cases}$

\item $\begin{cases}7x_1-x_2=6\\x_1+3x_2=8\end{cases}$

\item $\begin{cases}3x_1+2x_2=11\\x_1-4x_2=-9\end{cases}$

\item $\begin{cases}-2x_1-3x_2=-7\\x_1+5x_2=12\end{cases}$

\item $\begin{cases}3x_1-x_2=5\\-x_1+4x_2=6\end{cases}$

\item $\begin{cases}x_1+4x_2=9\\-3x_1+2x_2=7\end{cases}$

\item $\begin{cases}8x_1-3x_2=5\\2x_1+x_2=7\end{cases}$

\item $\begin{cases}2x_1-3x_2=-4\\-x_1+x_2=1\end{cases}$

\end{enumerate}
\end{multicols}

\emph{2. Solve the following $3\times 3$ systems by the matrix method.}
\bigskip
\begin{multicols}{2}\raggedright
\begin{enumerate}[label=\arabic*., leftmargin=*]

\item $\begin{cases}x_1-x_2+x_3=0\\2x_2-x_3=1\\2x_1+3x_2=5\end{cases}$

\item $\begin{cases}2x_1-x_2+3x_3=9\\x_1-x_3=-2\\3x_1+2x_2+x_3=7\end{cases}$

\item $\begin{cases}x_1+2x_2-x_3=3\\2x_1-x_2+3x_3=1\\x_1+x_2+x_3=6\end{cases}$

\item $\begin{cases}-x_1+x_2+x_3=1\\x_1-x_2+x_3=3\\x_1+x_2-x_3=5\end{cases}$

\item $\begin{cases}3x_1-x_2+2x_3=5\\x_1+2x_2-x_3=0\\2x_1+x_2+3x_3=11\end{cases}$

\item $\begin{cases}x_1+x_2+x_3=6\\x_1-x_2+2x_3=5\\2x_1+x_2-x_3=1\end{cases}$

\item $\begin{cases}2x_1+3x_2-x_3=4\\x_1-x_2+2x_3=7\\3x_1+x_2+x_3=10\end{cases}$

\item $\begin{cases}4x_1-x_2+2x_3=8\\x_1+3x_2-x_3=5\\-2x_1+x_2+3x_3=1\end{cases}$

\item $\begin{cases}x_1-2x_2+3x_3=7\\-x_1+x_2-x_3=-3\\2x_1+3x_2-2x_3=-1\end{cases}$

\item $\begin{cases}3x_1+2x_2+x_3=10\\x_1-x_2+4x_3=12\\2x_1+x_2-3x_3=-5\end{cases}$

\item $\begin{cases}x_1+3x_2+2x_3=8\\2x_1-x_2+x_3=3\\-x_1+2x_2-3x_3=-7\end{cases}$

\item $\begin{cases}x_1-x_2+x_3=2\\2x_1+x_2-3x_3=-1\\x_1+2x_2+x_3=7\end{cases}$

\item $\begin{cases}5x_1-2x_2+x_3=3\\x_1+3x_2-2x_3=1\\-x_1+x_2+4x_3=8\end{cases}$

\item $\begin{cases}-2x_1+x_2+3x_3=4\\x_1+2x_2-x_3=5\\3x_1-x_2+2x_3=6\end{cases}$

\item $\begin{cases}x_1+x_2-2x_3=-3\\3x_1-2x_2+x_3=4\\2x_1+x_2+3x_3=11\end{cases}$

\item $\begin{cases}2x_1-3x_2+x_3=0\\x_1+x_2-2x_3=-4\\-x_1+2x_2+x_3=5\end{cases}$

\item $\begin{cases}x_1+4x_2-x_3=7\\-2x_1+x_2+3x_3=2\\x_1-x_2+2x_3=3\end{cases}$

\item $\begin{cases}3x_1-x_2-2x_3=-1\\x_1+2x_2+x_3=8\\-x_1+3x_2-x_3=5\end{cases}$

\item $\begin{cases}-x_1-2x_2+x_3=1\\2x_1+x_2-3x_3=0\\-x_1-x_2+2x_3=3\end{cases}$

\item $\begin{cases}4x_1+x_2-x_3=6\\x_1-3x_2+2x_3=-1\\2x_1+2x_2+x_3=9\end{cases}$

\item $\begin{cases}x_1+x_2+x_3=9\\2x_1+x_2-x_3=2\\x_1-2x_2+3x_3=13\end{cases}$

\item $\begin{cases}3x_1+x_2-2x_3=3\\-x_1+2x_2+x_3=4\\x_1+x_2+3x_3=10\end{cases}$

\item $\begin{cases}2x_1-x_2+x_3=4\\x_1+3x_2-2x_3=-1\\-x_1+x_2+x_3=2\end{cases}$

\item $\begin{cases}-x_1+x_2+4x_3=5\\2x_1+3x_2-x_3=7\\x_1-2x_2+x_3=0\end{cases}$

\item $\begin{cases}x_1-3x_2+2x_3=1\\3x_1+x_2-x_3=8\\2x_1-x_2+3x_3=7\end{cases}$

\item $\begin{cases}x_1+2x_2-3x_3=-4\\-2x_1+x_2+x_3=3\\3x_1-x_2+2x_3=5\end{cases}$

\item $\begin{cases}2x_1+x_2+2x_3=10\\2x_1-x_2+x_3=3\\-x_1+2x_2-x_3=0\end{cases}$

\item $\begin{cases}x_1-x_2-x_3=-3\\-x_1+2x_2+3x_3=10\\2x_1+x_2-2x_3=-1\end{cases}$

\item $\begin{cases}3x_1+2x_2-x_3=7\\x_1-x_2+3x_3=6\\-2x_1+x_2+x_3=1\end{cases}$

\item $\begin{cases}x_1+x_2-x_3=4\\-x_1+2x_2+x_3=1\\2x_1-x_2+3x_3=5\end{cases}$

\end{enumerate}
\end{multicols}
\emph{3. Solve the systems $Ax=b$ with the given matrix $A$ for each of the given vectors $b$, computing $A^{-1}$ only once.}
\bigskip
\begin{enumerate}[label=\arabic*., leftmargin=*]

\item $A=\begin{pmatrix}2&1\\1&1\end{pmatrix}$; \quad
(a) $b=\begin{pmatrix}3\\2\end{pmatrix}$, \quad
(b) $b=\begin{pmatrix}5\\4\end{pmatrix}$, \quad
(c) $b=\begin{pmatrix}-1\\3\end{pmatrix}$.

\item $A=\begin{pmatrix}1&-1\\3&2\end{pmatrix}$; \quad
(a) $b=\begin{pmatrix}4\\1\end{pmatrix}$, \quad
(b) $b=\begin{pmatrix}0\\5\end{pmatrix}$, \quad
(c) $b=\begin{pmatrix}-3\\7\end{pmatrix}$.

\item $A=\begin{pmatrix}5&2\\-1&3\end{pmatrix}$; \quad
(a) $b=\begin{pmatrix}7\\4\end{pmatrix}$, \quad
(b) $b=\begin{pmatrix}11\\-2\end{pmatrix}$, \quad
(c) $b=\begin{pmatrix}0\\-1\end{pmatrix}$.

\item $A=\begin{pmatrix}3&-1\\4&1\end{pmatrix}$; \quad
(a) $b=\begin{pmatrix}2\\9\end{pmatrix}$, \quad
(b) $b=\begin{pmatrix}-1\\3\end{pmatrix}$, \quad
(c) $b=\begin{pmatrix}8\\0\end{pmatrix}$.

\item $A=\begin{pmatrix}-2&1\\1&3\end{pmatrix}$; \quad
(a) $b=\begin{pmatrix}0\\7\end{pmatrix}$, \quad
(b) $b=\begin{pmatrix}5\\-1\end{pmatrix}$, \quad
(c) $b=\begin{pmatrix}-4\\2\end{pmatrix}$.

\item $A=\begin{pmatrix}1&-1&1\\0&2&-1\\2&3&0\end{pmatrix}$; \quad
(a) $b=\begin{pmatrix}0\\1\\5\end{pmatrix}$, \quad
(b) $b=\begin{pmatrix}2\\-3\\1\end{pmatrix}$.

\item $A=\begin{pmatrix}1&2&0\\0&1&-3\\-1&0&2\end{pmatrix}$; \quad
(a) $b=\begin{pmatrix}1\\4\\-2\end{pmatrix}$, \quad
(b) $b=\begin{pmatrix}3\\0\\5\end{pmatrix}$.

\item $A=\begin{pmatrix}3&0&1\\-1&2&0\\0&-4&3\end{pmatrix}$; \quad
(a) $b=\begin{pmatrix}4\\-1\\7\end{pmatrix}$, \quad
(b) $b=\begin{pmatrix}-2\\3\\0\end{pmatrix}$.

\item $A=\begin{pmatrix}4&-1\\2&3\end{pmatrix}$; \quad
(a) $b=\begin{pmatrix}6\\8\end{pmatrix}$, \quad
(b) $b=\begin{pmatrix}-3\\11\end{pmatrix}$, \quad
(c) $b=\begin{pmatrix}1\\-5\end{pmatrix}$.

\item $A=\begin{pmatrix}1&0&-2\\0&3&1\\2&-1&0\end{pmatrix}$; \quad
(a) $b=\begin{pmatrix}5\\-2\\1\end{pmatrix}$, \quad
(b) $b=\begin{pmatrix}0\\4\\-3\end{pmatrix}$.

\item $A=\begin{pmatrix}7&-2\\3&-1\end{pmatrix}$; \quad
(a) $b=\begin{pmatrix}3\\1\end{pmatrix}$, \quad
(b) $b=\begin{pmatrix}-5\\0\end{pmatrix}$, \quad
(c) $b=\begin{pmatrix}9\\4\end{pmatrix}$.

\item $A=\begin{pmatrix}-1&2\\3&-5\end{pmatrix}$; \quad
(a) $b=\begin{pmatrix}1\\-4\end{pmatrix}$, \quad
(b) $b=\begin{pmatrix}7\\0\end{pmatrix}$, \quad
(c) $b=\begin{pmatrix}-2\\6\end{pmatrix}$.

\item $A=\begin{pmatrix}2&-1&0\\1&0&-1\\3&2&1\end{pmatrix}$; \quad
(a) $b=\begin{pmatrix}1\\-1\\4\end{pmatrix}$, \quad
(b) $b=\begin{pmatrix}3\\0\\-2\end{pmatrix}$.

\item $A=\begin{pmatrix}8&5\\3&2\end{pmatrix}$; \quad
(a) $b=\begin{pmatrix}1\\0\end{pmatrix}$, \quad
(b) $b=\begin{pmatrix}13\\5\end{pmatrix}$, \quad
(c) $b=\begin{pmatrix}-7\\-3\end{pmatrix}$.

\item $A=\begin{pmatrix}-2&1&3\\4&0&-1\\1&-2&5\end{pmatrix}$; \quad
(a) $b=\begin{pmatrix}0\\3\\-1\end{pmatrix}$, \quad
(b) $b=\begin{pmatrix}5\\-2\\4\end{pmatrix}$.

\item $A=\begin{pmatrix}1&1\\1&2\end{pmatrix}$; \quad
(a) $b=\begin{pmatrix}3\\5\end{pmatrix}$, \quad
(b) $b=\begin{pmatrix}0\\-1\end{pmatrix}$, \quad
(c) $b=\begin{pmatrix}10\\15\end{pmatrix}$.

\item $A=\begin{pmatrix}4&-1&2\\3&5&0\\-2&0&1\end{pmatrix}$; \quad
(a) $b=\begin{pmatrix}7\\-1\\3\end{pmatrix}$, \quad
(b) $b=\begin{pmatrix}0\\8\\-4\end{pmatrix}$.

\item $A=\begin{pmatrix}3&4\\2&3\end{pmatrix}$; \quad
(a) $b=\begin{pmatrix}1\\0\end{pmatrix}$, \quad
(b) $b=\begin{pmatrix}7\\5\end{pmatrix}$, \quad
(c) $b=\begin{pmatrix}-2\\3\end{pmatrix}$.

\item $A=\begin{pmatrix}0&-3&1\\5&2&0\\-4&0&3\end{pmatrix}$; \quad
(a) $b=\begin{pmatrix}2\\5\\-1\end{pmatrix}$, \quad
(b) $b=\begin{pmatrix}-3\\0\\7\end{pmatrix}$.

\item $A=\begin{pmatrix}6&1\\-2&3\end{pmatrix}$; \quad
(a) $b=\begin{pmatrix}7\\1\end{pmatrix}$, \quad
(b) $b=\begin{pmatrix}13\\-5\end{pmatrix}$, \quad
(c) $b=\begin{pmatrix}0\\6\end{pmatrix}$.

\item $A=\begin{pmatrix}1&3&-1\\-2&0&4\\0&5&2\end{pmatrix}$; \quad
(a) $b=\begin{pmatrix}4\\-6\\3\end{pmatrix}$, \quad
(b) $b=\begin{pmatrix}-1\\2\\7\end{pmatrix}$.

\item $A=\begin{pmatrix}5&-2\\-3&1\end{pmatrix}$; \quad
(a) $b=\begin{pmatrix}1\\0\end{pmatrix}$, \quad
(b) $b=\begin{pmatrix}-4\\3\end{pmatrix}$, \quad
(c) $b=\begin{pmatrix}11\\-7\end{pmatrix}$.

\item $A=\begin{pmatrix}3&1&0\\-2&4&-1\\0&-3&5\end{pmatrix}$; \quad
(a) $b=\begin{pmatrix}1\\6\\-2\end{pmatrix}$, \quad
(b) $b=\begin{pmatrix}-5\\0\\3\end{pmatrix}$.

\item $A=\begin{pmatrix}0&1\\-1&3\end{pmatrix}$; \quad
(a) $b=\begin{pmatrix}2\\4\end{pmatrix}$, \quad
(b) $b=\begin{pmatrix}-1\\0\end{pmatrix}$, \quad
(c) $b=\begin{pmatrix}5\\-3\end{pmatrix}$.

\item $A=\begin{pmatrix}-1&2&-3\\4&0&1\\0&-5&2\end{pmatrix}$; \quad
(a) $b=\begin{pmatrix}0\\9\\-5\end{pmatrix}$, \quad
(b) $b=\begin{pmatrix}3\\-1\\2\end{pmatrix}$.

\item $A=\begin{pmatrix}9&-1\\4&1\end{pmatrix}$; \quad
(a) $b=\begin{pmatrix}8\\5\end{pmatrix}$, \quad
(b) $b=\begin{pmatrix}0\\-3\end{pmatrix}$, \quad
(c) $b=\begin{pmatrix}17\\9\end{pmatrix}$.

\item $A=\begin{pmatrix}2&0&-1\\0&3&4\\-3&1&0\end{pmatrix}$; \quad
(a) $b=\begin{pmatrix}1\\-2\\4\end{pmatrix}$, \quad
(b) $b=\begin{pmatrix}-3\\5\\0\end{pmatrix}$.

\item $A=\begin{pmatrix}-3&2\\4&-3\end{pmatrix}$; \quad
(a) $b=\begin{pmatrix}1\\-1\end{pmatrix}$, \quad
(b) $b=\begin{pmatrix}7\\-8\end{pmatrix}$, \quad
(c) $b=\begin{pmatrix}-5\\4\end{pmatrix}$.

\item $A=\begin{pmatrix}1&-4&0\\3&2&-5\\-2&0&1\end{pmatrix}$; \quad
(a) $b=\begin{pmatrix}2\\-3\\1\end{pmatrix}$, \quad
(b) $b=\begin{pmatrix}-1\\6\\0\end{pmatrix}$.

\item $A=\begin{pmatrix}2&-1\\-1&1\end{pmatrix}$; \quad
(a) $b=\begin{pmatrix}3\\0\end{pmatrix}$, \quad
(b) $b=\begin{pmatrix}-2\\5\end{pmatrix}$, \quad
(c) $b=\begin{pmatrix}7\\-4\end{pmatrix}$.

\end{enumerate}
 